\begin{enumerate}
\item \textbf{Quelques exemples.}
\begin{enumerate}
    \item $1 + 2 + \cdots + 5 = 15$ et $7 + 8 = 15$ donc 6 est équilibré de balance 2.
    
    \item $1 + 2 + \cdots + 6 = 21$ et $8 + 9 = 17 < 21 < 27 = 8 + 9 + 10$.
    
    \item On a
    \[1 + 2 + 3 + \cdots + 203 = \frac{203 \times (203 + 1)}{2} = 20706\]
    et
    \[(204 + 1) + (204 + 2) + \cdots + (204 + 84) = \frac{84 \times (2 \times 204 + 84 + 1)}{2} = 20416\]
\end{enumerate}

\item \textbf{}
\begin{enumerate}
    \item Si $n$ est équilibré de balance $b$, on a :
    \[\frac{(n - 1)n}{2} = \frac{b(2n + b + 1)}{2}\]
    donc
    \[n^2 - n = 2bn + b^2 + b\]
    On en déduit que
    \[b^2 + (2n + 1)b + (n - n^2) = 0\]
    donc $b$ est solution d'une équation du second degré dont le discriminant est égal à $8n^2 + 1$ qui est strictement positif. L'équation admet donc deux solutions :
    \[\frac{-(2n + 1) + \sqrt{8n^2 + 1}}{2} \quad \text{et} \quad \frac{-(2n + 1) - \sqrt{8n^2 + 1}}{2}\]
    Comme $b$ est un entier strictement positif, on a nécessairement $b = \dfrac{-(2n+1)+\sqrt{8n^2+1}}{2}$.
    
    On en déduit que $8n^2 + 1 = (2b + 2n + 1)^2$ donc $8n^2 + 1$ est un carré parfait.
    
    \item Comme $e^2 = 8n^2 + 1$, $e^2$ est impair, ce qui impose que $e$ l'est aussi. On en déduit que $b = \dfrac{e - (2n + 1)}{2}$ est un entier. Montrer que $b > 0$ revient à établir que $2n + 1 < \sqrt{8n^2 + 1}$, ce qui est aisé en comparant les carrés de ces deux nombres positifs (on rappelle que $n > 1$).
    
    \item La question 2.a donne le sens direct tandis que la question 2.b fournit la réciproque : puisque $b = \dfrac{-(2n+1)+\sqrt{8n^2+1}}{2}$ on a une solution de l'équation
    \[x^2 + (2n + 1)x + (n - n^2) = 0\]
    ce qui montre que
    \[b^2 + (2n + 1)b + (n - n^2) = 0\]
    On en déduit finalement que
    \[1 + 2 + 3 + \cdots + (n - 1) = \frac{(n - 1)n}{2} = \frac{b(2n + b + 1)}{2} = (n + 1) + (n + 2) + \cdots + (n + b)\]
    ce qui montre que $n$ est un nombre équilibré de balance $b$.
\end{enumerate}

\item \textbf{}
\begin{enumerate}
    \item On calcule :
    \begin{align*}
    8f(x)^2 + 1 &= 8\left(9x^2 + 6x\sqrt{8x^2 + 1} + 8x^2 + 1\right) + 1 \\
    &= 136x^2 + 48x\sqrt{8x^2 + 1} + 9
    \end{align*}
    et
    \begin{align*}
    (8x + 3\sqrt{8x^2 + 1})^2 &= 64x^2 + 48x\sqrt{8x^2 + 1} + 9(8x^2 + 1) \\
    &= 136x^2 + 48x\sqrt{8x^2 + 1} + 9
    \end{align*}
    
    \item Soit $n$ un nombre équilibré. D'après la question 2.a, $f(n)$ est bien un entier. On remarque par ailleurs que, comme $n \geq 2$, on a : $f(n) \geq 2$. Enfin, d'après la question 3.a, $8f(n)^2 + 1$ est un carré parfait. D'après la question 2.b, on conclut que $f(n)$ est un nombre équilibré.
    
    \item On sait que 6 est un nombre équilibré. Par ailleurs, pour tout entier naturel non nul $k$, $\dfrac{f(k)}{k} > 1$ donc $f(k) > k$. On en déduit que les nombres
    \[6, f(6), f(f(6)), f(f(f(6))), \ldots\]
    sont des nombres équilibrés (question 3.b) deux à deux distincts classés dans l'ordre croissant, ce qui donne le résultat.
    
    \item On sait que $u_k = 3u_{k-1} + \sqrt{8u_{k-1}^2 + 1}$ donc $(u_k - 3u_{k-1})^2 = 8u_{k-1}^2 + 1$. On obtient, après simplifications : $u_{k-1}^2 - 6u_k u_{k-1} + u_k^2 - 1 = 0$ donc $u_{k-1}$ est solution d'une équation du second degré dont le discriminant est égal à $4(8u_k^2 + 1)$ qui est strictement positif. L'équation donne donc deux solutions, qui sont
    \[3u_k + \sqrt{8u_k^2 + 1} \quad \text{et} \quad 3u_k - \sqrt{8u_k^2 + 1}\]
    Comme $u_{k-1} < u_k$, on a donc nécessairement : $u_{k-1} = 3u_k - \sqrt{8u_k^2 + 1}$.
    
    \item On remarque que $h$ est dérivable et que pour tout $x > 0$
    \[h'(x) = \frac{3}{2\sqrt{x}} - \frac{4}{\sqrt{8x + 1}}\]
    Or, $(3\sqrt{8x + 1})^2 = 9(8x + 1) = 72x + 9$, $(8\sqrt{x})^2 = 64x$ donc $3\sqrt{8x + 1} - 8\sqrt{x} > 0$. Comme $2\sqrt{x}\sqrt{8x + 1} > 0$, on en déduit que $h'(x) > 0$. Ainsi $h$ est strictement croissante. Pour tous réels $x$ et $y$ tels que $0 < x < y$, on a : $0 < x^2 < y^2$. Donc, comme $h$ est strictement croissante, $h(x^2) < h(y^2)$ donc $g(x) < g(y)$. Ainsi $g$ est strictement croissante.
    
    \item $u_2 = f(u_1) = f(6) = 35$. Puisque $n$ est équilibré, on déduit du résultat admis par l'énoncé que si $n \leq 35$, alors $n = 6$ ou $n = 35$. Or, $u_1 = 6$ et $u_2 = 35$ et $n$ n'est pas de la forme $u_k$. On en déduit que $n > 35$, c'est-à-dire que $n > u_2$. Puisque la suite d'entiers naturels $(u_k)_{k \geq 1}$ est strictement croissante, il existe un entier $m \geq 2$ tel que $u_m < n < u_{m+1}$.
    
    On déduit de la croissance de $g$ établie à la question 3.e que $g(u_m) < g(n) < g(u_{m+1})$. Or, d'après la question 3.d, $g(u_m) = u_{m-1}$ et $g(u_{m+1}) = u_m$. D'où $u_{m-1} < g(n) < u_m$. Comme $u_m < n$, on a donc $g(n) < u_m < n$ donc $g(n) < n$.
    
    D'après la question 3.b, $n > 35$ donc $n > 6$. Or, $n$ est un nombre équilibré. On vérifie (comme on l'a fait avec $f(n)$) que $g(n)$ l'est aussi. On en déduit que $g(n)$ est un nombre équilibré qui ne s'écrit pas sous la forme $u_k$ pour un certain entier $k \geq 1$. Finalement $n$ n'est pas le plus petit nombre équilibré ne s'écrivant pas sous forme $u_k$ pour un certain entier $k \geq 1$, ce qui est absurde par définition de $n$. Ainsi, tout nombre équilibré est de la forme $u_k$ pour un certain entier $k \geq 1$.
\end{enumerate}

\item \textbf{}
\begin{enumerate}
    \item Soit $k \geq 2$ un entier. Alors $u_{k+1} = 3u_k + \sqrt{8u_k^2 + 1}$ et, d'après la question 3.d, $u_{k-1} = 3u_k - \sqrt{8u_k^2 + 1}$. Donc $u_{k+1} = 3u_k + (3u_k - u_{k-1}) = 6u_k - u_{k-1}$.
    
    \item Dans $s$, on calcule $s = 1 + \ldots + i$ jusqu'à ce que $i$ arrive à $n$ (exclu). La fonction \texttt{mystere(n)} calcule donc la somme $1 + 2 + 3 + \ldots + (n - 1)$.
    
    \item On déduit de la question 4.b la fonction suivante.
    
    \begin{verbatim}
def equilibre(n):
    k = mystere(n)   # k vaut 1+2+...+(n-1)
    s = n+1
    i = n+2
    while s < k:
        s = s + i
        i = i + 1
    return s == k
    \end{verbatim}
    
    Dans $s$, on calcule $s = (n+1) + \ldots + i$ jusqu'à ce que cette somme dépasse le seuil $k$. On regarde alors si la balance est équilibrée.
    
    Mais on peut aussi écrire la fonction suivante, de complexité bien meilleure, en exploitant la question 4.a.
    
    \begin{verbatim}
def equilibre(n):
    (a, b) = (6, 35)
    while a < n:
        (a, b) = (b, 6*b - a)
    return a == n
    \end{verbatim}
\end{enumerate}

\item \textbf{} On propose ici aussi deux solutions, la seconde étant plus efficace que la première.

La première fonction utilise la question 4.c.

\begin{verbatim}
def liste_equilibres(n):
    l = []
    for i in range(2, n):
        if equilibre(i):
            l.append(i)
    return l
\end{verbatim}

La seconde fonction exploite encore la question 4.a.

\begin{verbatim}
def liste_equilibres(n):
    l = []
    (a, b) = (6, 35)
    while a < n:
        l.append(a)
        (a, b) = (b, 6*b - a)
    return l
\end{verbatim}

\textbf{Remarques}

La balance de 6930 est 2870.

La balance de 235416 est 97512.

La balance de 313506783024 est 129858761424.

\texttt{liste\_equilibres(1000000000000)} renvoie

\texttt{[6, 35, 204, 1189, 6930, 40391, 235416, 1372105, 7997214, 46611179, 271669860, \\1583407981, 9228778026, 53789260175, 313506783024]}




\end{enumerate}
